% Generated from paper/source. Do not edit by hand.
\section{Introduction}
Claims about a new neural representation are only as strong as the control that could have produced the same computation. This is especially important for complex-valued models. Their states can express magnitude and phase compactly, and complex multiplication naturally represents rotations and ordered transformations. Yet a complex linear map also has an exact structured representation over twice as many real coordinates. A comparison against an ordinary real recurrent cell, or against two real channels that are not coupled by the corresponding block matrix, cannot identify whether any observed advantage comes from complex arithmetic, from a structured constraint, from optimization, or simply from an underpowered control.

Q-Neuro began with a positive result. In a family of synthetic sequential reasoning environments called NeuroWorld, an ordered complex operator was more robust under declared causal shifts than several then-tested controls. At moderate shift the mean complex-minus-two-channel-real difference was +0.0602 top-1 accuracy across five worlds. The result replicated within that simulator, and destructive phase interventions showed that the trained model used its phase degrees of freedom. Neither observation established that complex arithmetic was necessary. Both remained compatible with a real structured implementation and with stronger real sequence models.

We therefore reorganized the project around a falsification-first question: \textbf{does the complex operator retain a practically relevant advantage when compared with an exact real implementation and a cellwise best-real envelope on task families not used to formulate the claim?} The answer within the executed scope is no. The exact real block matches the complex computation at numerical precision. Once it is admitted to the comparator set, the sign of the performance gap reverses in a power pilot, in reduced discovery, and in held-out reduced confirmation. A quantitative candidate law that appeared accurate in discovery fails its frozen held-out thresholds. Finally, the preregistered grand benchmark is not opened because the program has not completed the full comparator, search-budget, compute-accounting, shift-grid, and raw-prediction requirements.

This paper makes four bounded contributions. First, it gives an implementation-level exact-real falsifier for the Q-Neuro operator and verifies mapped equivalence across independent task families. Second, it constructs eight nonmedical sequential generator families and a severity-indexed shift surface so that discovery and confirmation can be separated by generator family rather than by random examples alone. Third, it preserves a prospective law failure without refitting and a grand-experiment stop decision without weakening the gates. Fourth, it releases a machine-readable claim ledger that pairs every conclusion with counterevidence, scope, uncertainty, and a prospective falsifier.

The work does not claim clinical validity, universal superiority of real networks, or evidence for a physical quantum process in cognition. The independent studies are deliberately reduced relative to the full preregistration: only five model families were trained, search used one configuration per model, training sizes were restricted, per-trial FLOPs and optimizer steps were not preserved, and raw predictions were not stored. These limitations block the grand experiment and external-validity claims. They do not alter the exact algebraic control or the observed absence of a complex advantage in the executed cells. The resulting paper is a case study in how a project can become scientifically stronger when its favored hypothesis fails.
